\chapter{Conclusion}
\label{chap:conclusion}

This chapter concludes the dissertation. 
We have taken the time to look at the problem of delays in \gls{rl}.
The results of this dissertation particularly highlight the relations between \glspl{dmdp} and \glspl{mdp}, \glspl{pomdp}, higher-order \glspl{mdp} or higher-order \glspl{mc}.
Our principal contributions, gathered along four axis, are summarised in the ensuing section. 
To enhance reader comprehension, these contributions will be linked to the section they originated from.
We end with a discussion of the potential future research directions.

\section{Discussion and Key Results}
\subsubsection{Unifying Framework for Delays}
The first contribution of this thesis was to provide a unified framework to the problem of delays.
With this objective in mind, in \Cref{chap:delay}, we proposed a versatile mathematical framework to bring together the different delayed processes encountered in the literature. 
This framework, called \gls{dmdp}, is based on an underlying \gls{mdp} to which a process is added that defines the delay.
This definition encompasses constant, non-integer, stochastic, or Markovian delays as particular cases. 
Then, we presented the main types of delay encountered in the literature and in practice (\Cref{sec:nature_delays}) and organised them according to the variable they affect--the state, the action or the reward--and the process that they follow--constant, stochastic, state-dependent processes...
Finally, we presented an in-depth discussion of delays in \gls{rl} and related fields such as real-time \gls{rl} and bandits, with a particular focus on state observation and action execution delays in \Cref{sec:related_works}.
For the latter, we gathered algorithms from the literature into three categories: the augmented state, the memoryless, and the model-based approaches.

\subsubsection{Theoretical Understanding of Delays}
Our first theoretical result on the effect of delayed state observation or action execution on the \gls{mdp}, is a demonstration of a central property of \gls{dmdp} with constant delay: the longer the delay, the lower the performance (\Cref{th:perf_delay_vector_all_weight}).
In addition, we showed that a longer constant delay implied less variance in the return or the average reward (\Cref{pp:var_ismmdp}). 
The articulation of these two results is well illustrated in \Cref{fig:evolution_delay}.

In \Cref{chap:multi_action_delay}, we explored a more general framework, which we call the multiple action delay.
This framework leverages the \gls{mtd} and \gls{mtdg} models from the \gls{mc} literature and applies them to \glspl{mdp} to obtain what we call \glspl{mtdmdp}.
We defined four such \glspl{mtdmdp}: \glspl{ismmdp}, \glspl{immmdp}, \glspl{psmmdp} and \glspl{pmmmdp} which differ in the way the \gls{mtd} and \gls{mtdg} models are applied.
We first showed that state augmentation can cast these problems back to an \gls{mdp} (\Cref{prop:equ_mdp_im_mtdmdp} and the following propositions).
Then, we focused on learning policies to optimise the expected discounted return and average rewards in these models. 
For \glspl{psmmdp} and \glspl{pmmmdp}, we showed that the delay was essentially ineffective (\Cref{th:psmmdp_same} and \Cref{th:pmmmdp_same}). 
In fact, an undelayed policy yields the same average reward, whether applied to these processes or to the underlying undelayed \gls{mdp}.
Consequently, we studied \glspl{ismmdp} which are more interesting, as they include constant delay \gls{dmdp} as a special case. 
Analysing the effect of the delay vector $\boldsymbol\lambda$ on these models, we conjectured that the cumulative delay distribution was the key to ordering the potential optimal return or average reward (\Cref{conj:delay_cumulativ_distrib}).  
We provided arguments for this claim and showed that the mean reward could not substitute the cumulative delay distribution in the conjecture.
To conclude the theoretical analysis, we examined how some properties of the underlying \gls{mdp} could be handed over to the \gls{ismmdp} built upon it.
We showed that a unichain \gls{mdp} did not imply a unichain \gls{ismmdp}(\Cref{pp:unichain_ismmdp}) while a communicating \gls{mdp} did imply a communicating \gls{ismmdp} (\Cref{pp:communicating_ismmdp}).
Leveraging these results, we conclude that UCRL cannot be readily applied to \gls{ismmdp}, while UCRL2 can.  
An empirical analysis concluded this chapter, highlighting the implications of theory and studying the effect of $\boldsymbol\lambda$ on both the expected discounted return and the average reward. 

Shifting the focus toward delayed \gls{rl} algorithms, we showed that model-based approaches can yield sub-optimal policies (\Cref{pp:counter_example_belief_based}).
Nevertheless, in \Cref{th:perf_diff_bound}, we provided lower-bounds performance guarantees for model-based approaches in smooth \glspl{mdp}, including for non-integer delays.
Notably, these bounds match the theoretical upper bound up to a multiplicative factor (\Cref{th:lb_tight}).
These results imply that the delay has at worst only an additive impact on performance loss, which is in line with the results from the bandit literature. 
Finally, we extended the guarantees to the case in which the model-based policy was trained on constant delays but tested on stochastic delays, producing the first performance bound in \gls{rl} for anonymous action execution delays (\Cref{thm:stoch_mdp_dida_bound}).


\subsubsection{Algorithms for Delayed Reinforcement Learning}
In \Cref{chap:belief_based}, we introduced a new model-based approach, the belief representation network.
Its idea is to learn a vectorial representation of the belief in two steps.
First, a Transformer network processes the augmented state and produces a vectorial representation of the belief for each future state up to the undelayed one.
Second, a masked-autoregressive flow network uses this representation to fit the distribution of the future states in order to ensure that the representation actually encodes the belief. 
Compared to previous model-based approaches that learn some statistics of the current unobserved state, knowledge of the belief contains more information and enables more complex policies. 
This belief representation can be readily plugged into any \gls{rl} algorithm as a replacement for the state, as we have done with \gls{trpo} for our approach called D-TRPO.
This idea was shown to be effective in experiments and demonstrates that many delays can be learnt in a single training, thus reducing the cost of taking into account delays (\Cref{subsec:results_belief}).
Further experiments have highlighted the ability of D-TRPO to adapt to a wide range of problems, particularly stochastic \glspl{dmdp}.
Interestingly, the experiments showed that the A-SAC baseline--an augmented approach--is very efficient in tackling delay (\Cref{subsec:results_belief}).
Finally, we have shown that the policy learnt by D-TRPO for a certain delay can be easily used to apply on smaller delays (\Cref{subsec:multi_delay_once_belief}).

In \Cref{chap:imitation_undelayed} we designed a simple yet effective algorithm, DIDA, which imitates an undelayed expert given the knowledge of the augmented state.
Using the imitation learning literature, we showed that \textsc{DAgger} was a suitable algorithm to perform the imitation step.
Although DIDA does not have access to the current state and cannot always perfectly mimic the expert, \Cref{th:perf_diff_bound} provides lower-bound performance guarantees for DIDA in smooth \glspl{mdp} with constant delay.
This result also applies to non-integer delays and \Cref{thm:stoch_mdp_dida_bound} provides guarantees for DIDA on stochastic delays.
This means that DIDA is suitable for a wide range of applications.
This versatility was confirmed by our extensive empirical study.
In a wide range of tasks, DIDA achieved state-of-the-art performance among a large number of baselines, including A-SAC.
Notably, DIDA is more sample efficient than these baselines and is not computationally expensive.
However, DIDA requires the knowledge or training of an undelayed expert, which could be a limitation in some problems.

In \Cref{chap:lifelong}, we explored the solution of memoryless policies.
Because the agent is blind to part of the state, to it, the dynamics are inherently non-stationary.
From this observation, we designed an algorithm that can adapt to non-stationary dynamics.
In particular, the non-stationarity arises intra-episode which makes the problem more similar to the lifelong setting.
The idea of the algorithm, named POLIS, is to handle non-stationarity at a hyper-policy level. 
The hyper-policy takes time as input and outputs the parameters of a policy to be queried at that time. 
In this way, policies remain stationary while the non-stationarity is taken into account inside the hyper-policy.
To optimise its hyper-policy, POLIS builds an estimator of the future by leveraging past data through multiple importance sampling, as explained in \Cref{subsec:mis_lifelong}.
In smooth environments, the estimator bias is bounded and can be controlled by a parameter $\omega$ (\Cref{lem:bias_bound}).
In addition to the estimated future return, two terms were included in the objective.
First, to prevent catastrophic forgetting, the past performance of the hyper-policy was added to the objective.
In this way, the agent is incentivised to keep a memory of past efficient behaviours.
Second, to avoid overfitting the past by learning an over non-stationary hyper-policy, we regularise the objective with a probabilistic upper-bound of the variance of the estimator (\Cref{th:bound_surrogate_objective}).
Finally, in \Cref{subsec:extension_delay_polis}, we detail how POLIS can be applied to a delayed process with simple modifications to the estimator.
Experimental results demonstrated that POLIS can adapt to non-stationary dynamics and avoid excessive non-stationarity when not necessary.
Applied to a delayed environment, POLIS shows robust performance that naturally decreases with longer delays.

\subsubsection{Extensive Empirical Evaluation}
Our empirical evaluation, conducted on a wide range of tasks and delays, gave useful information on the peculiarities and abilities of delayed \gls{rl} algorithms.
These tasks included deterministic environments such as classic control in the Pendulum environment, pathfinding in a maze, or robotic locomotion in Mujoco (\Cref{chap:belief_based}, \Cref{chap:imitation_undelayed}, \Cref{chap:multi_action_delay}).
They also included different adaptations of Pendulum, with stochastic delays (\Cref{chap:belief_based}, \Cref{chap:imitation_undelayed}), non-integer delays (\Cref{chap:imitation_undelayed}) or multiple action delays (\Cref{chap:multi_action_delay}).
More realistic environments were also considered, such as FOREX trading (\Cref{chap:imitation_undelayed},\Cref{chap:lifelong}) and water resource management (\Cref{chap:lifelong}).

On these tasks, we studied our algorithms as well as many baselines from the literature, including augmented state \gls{trpo} (A-TRPO), memoryless \gls{trpo} (M-TRPO), augmented state \gls{sac} (A-SAC),  memoryless \gls{sac} (M-SAC), SARSA, dSARSA, FQI, Pro-WLS, LPG-FTW and UCRL2.

\section{Future Research Directions}
Before concluding this dissertation, we recall some of the future directions that we have proposed in the previous chapters. 
For model-based approaches, such as D-TRPO, a valuable future direction could be to leverage the environment model to plan for a longer horizon.
For instance, this could be used to enhance the value function estimation in actor-critic methods.
Regarding DIDA, the strongest limitation is the need for an undelayed expert. 
Potential future work could consider learning an undelayed expert offline, from a dataset collected by a delayed policy. 
In this way, the whole interaction with the environment would always be made in the delayed case. 
For POLIS, a further source of non-stationary could be added by considering stochastic delays.
The advantage of a memoryless policy in this case is that its input is constant and ``small'' in dimension $\cardinal{\statespace}$.
Finally, for \glspl{mtdmdp}, a future directions could be to leverage techniques for the estimation of the parameters of an \gls{mtd} model \cite{berchtold2002mixture} or the reward structure \cite{romano2022multi} to enhance theoretical \gls{rl} guarantees. 

I hope that this dissertation has been useful in providing a global overview of state observation and action execution delays and in broadening their understanding. 
I am excited to see the next developments of \gls{rl}, both theoretically and in its applications, and, of course, I will be particularly interested in the way the delay will be dealt with in these developments.

\section{Final Word}
Given the substantial practical importance of delays on performance and the existence of straightforward mechanisms to address them, I strongly advocate that readers integrate delay considerations into their \gls{rl} applications.


%=============================================
% For notations
%---------------------------------------------

\begin{align*}
    \vphantom{
        \text{
            \gls{delay}
        }
    }
\end{align*}
\begin{align*}
    \vphantom{
        \text{
            \gls{augmentedstatespace}
        }
    }
    \vphantom{
        \text{
            \gls{actionspace}
        }
    }
    \vphantom{
        \text{
            \gls{statespace}
        }
    }
\end{align*}

